\cotefile{161-2}

\scanpage{4}

\uline{Théorème} Soient $d=(d',d'')$ un \inserted{couple de} petits ens de diagrammes, $C$ une catégorie ens. petite $\delta=(\delta',\delta'')$ un \inserted{petit} ens de cônes projectifs et inductifs dans $C$ $\Delta=\Delta_d=\Delta_{d',d''}$ $2$-Catégorie des \unsure{fatigues|catégories} ayant des lim de type $d$, avec foncteurs y commutant

\deleted{Alors il existe}

Pour tte $E\in\mathrm{Ob}\,\Delta$, soit $\underline{\mathrm{Hom}}_{\deleted{\Delta}\delta}(C,E)$ la sous-cat. \uline{pleine} de $\underline{\mathrm{Hom}}_{\deleted{\unsure{\mathrm{Ens}}}}(C,E)$ formée des foncteurs qui sont $\delta$-exacts. Si $\bar C_{\deleted{\unsure{\delta}}}$ est de même munie de $\bar\delta=(\bar\delta',\bar\delta'')$ et est $\in\Delta$, on définit $\underline{\mathrm{Hom}}_{\Delta,\delta_1}(C_1,E)$ comme intersection de $\underline{\mathrm{Hom}}_{\bar\delta}$ et $\underline{\mathrm{Hom}}_\Delta$. Ceci posé, il existe une catégorie $\bar C_{\deleted{\delta}}\in\mathrm{Ob}\,\Delta$, et un foncteur $i:C\unsure{-|\to}\bar C_{\deleted{\delta}}$, tel que

a) $i$ soit $\delta$-exact i.e. $i\in\underline{\mathrm{Hom}}_\delta(C,\bar C_{\deleted{\delta}})$

b) il soit universel en le sens suivant : Pour tt $E\in\mathrm{Ob}\,\Delta$, le foncteur $\varphi\longmapsto\varphi\circ i$ induit une équivalence

\[
\underline{\mathrm{Hom}}_{\Delta,i(\delta)}(\bar C_{\deleted{\delta}},E)\longrightarrow\underline{\mathrm{Hom}}_\delta(C,E)
\]

[ i.e. $(\bar C,i)$ est objet $2$-initial de la $2$-catégorie des $E$ \inserted{$\in\mathrm{Ob}\,\Delta$} sous $C$ tels que $C\to E$ soit $\delta$-exact ]. De plus, $\bar C_{\deleted{\delta}}$ est ens. petite.

\uline{NB} Ce Pb résoud à la fois le \uline{Pb 1} $(\unsure{\Gamma'|\delta'}=\delta''=\varnothing)$ et le \uline{Pb} 2 ($C\in\mathrm{Ob}\,\Delta$, \deleted{\illegible{1}} $\delta'=$ \deleted{diagrammes} \inserted{cônes pr. de type $d'$} exacts \deleted{\unsure{égaux à|égaux}} $\cup$ les diagrammes $(x\to y)\in\Sigma$, $\delta''=$ \deleted{diagrammes} cônes inductifs de type $d''$ exacts). On verra que le th. permet aussi de traiter le cas \deleted{\unsure{des|de}} conditions d'exactitude.

\uline{Dém} On construit \deleted{\unsure{une}} \inserted{par} récurrence transfinie \deleted{\unsure{et|une}} suite transfinie de \inserted{petites} catégories $C_i$, sous $C$, munies de $\delta_i=(\delta_i',\delta_i'')$ (petit).

\diagram{\begin{tikzcd}[ampersand replacement=\&] C\arrow[r,"\alpha_i"]\& C_i\arrow[r,"\tau_i",dash]\& C_{i+1}\end{tikzcd}}

\[
\left\{\begin{aligned}
&C_0=C,\ \delta_0=\delta=(\delta',\delta''),\\
&\deleted{(C_{i+1},\delta_{i+1})\text{ déduit de }(C_i,\delta_i)\text{ ainsi :}}\\
&(C_{i+1},\delta_{i+1})=K(C_i,\delta_i),\quad K\text{ étant explicité ci-dessous}\\
&C_{i\deleted{\unsure{+1}}}=\varinjlim_{j<i\,\mathrm{Cat}} C_i\quad\text{si }i\text{ ordinal limite},\quad
\begin{aligned}\delta_i'&=\bigcup\text{ images }\delta_j'\text{ dans } C_i\\
\delta_i''&=\bigcup\text{ im. }\delta_j''\text{ ds }C_i\end{aligned}
\end{aligned}\right.
\]

% archive sheet number: 3
\scanpage{5}

La propriété essentielle de $K(C,\delta)$ \deleted{\unsure{a}} est résumée \unsure{dans} :

a) $\tau_i:C_i\unsure{-|\to}C_{i+1}$ est \deleted{\unsure{préexact pour $\delta$|préexact}} \inserted{$\delta_i$-préexact}, i.e.

\[
a')\left\{\begin{aligned}
a'_1)&\ \deleted{\text{Unicité de ttes flèches }z\rightrightarrowsx}\ \text{Pour }(x\xleftarrow{p_\lambda}x_\lambda)\in d'_i,\text{ et ttes doubles}\\
&\text{flèches }z\overset{u,v}{\rightrightarrows}x\text{ telles que }p_\lambda\circ u=p_\lambda\circ v\quad\forall\lambda,\text{ on a }\tau_i(u)=\tau_i(v)\\
a'_2)&\ \text{Pour tte }\deleted{\unsure{\text{sys}}}\ \text{\uline{augmentation} }(z\xrightarrow{q_\lambda}x_\lambda)\text{ du syst. proj., il existe}\\
&\tau_i(z)\xrightarrow{q}\tau_i(x)\text{ telle que }\tau_i(p_\lambda)\circ q=\tau_i(q_\lambda)\quad\forall\lambda
\end{aligned}\right.
\]

a$''$) \deleted{\unsure{comme}} dual de a\unsure{′} ), avec $d''_i$ au lieu de $d'_i$.

\deleted{b) Tout diagramme $\unsure{d'_i|d'_{i+1}},\unsure{\Theta|\supset}\in d'\quad\delta'_{i+1}\supset\tau_i(\delta'_i)$}

\[
b)\left\{\begin{aligned}
b')&\ \forall\text{ diagramme }\inserted{\Theta'}\text{ de type }d'\text{ dans } C_i,\ \exists\text{ une }\deleted{\unsure{\text{augmen.}}}\inserted{\text{augmentation}}\text{ de}\\
&\tau_i(\Theta')\text{ qui est }\in\delta'_{i+1}\\
b'')&\ \forall\text{ diagramme }\Theta''\text{ de type }d''\text{ dans }C_i,\ \exists\text{ une augmentation de }\tau_i(\Theta'')\inserted{\text{ dans }\delta''_{i+1}}.
\end{aligned}\right.
\]

\[
c)\left\{\begin{aligned}
&\delta'_{i+1}\ \deleted{=}\ \inserted{\text{est}}\ \deleted{\unsure{\text{formé}}}\ \tau_i(\delta'_i)\cup\text{ diagr. }\deleted{\unsure{\in\delta'_{i+1}}}\text{ }\unsure{\text{munis}|\text{nouveaux}}\text{ de type }d'\text{ construits dans }b',\\
&\delta''_{i+1}=\tau_i(\delta''_i)\cup\text{ diagrammes de type }d''.\\
&\deleted{\text{tout élément de }\delta'_{i+1}\ (\delta''_{i+1})\text{ de type }\Theta'\ \unsure{\notin}d'\ (\Theta''\notin d'')\text{ est dans }\tau_i(\delta'_i),\ (\tau_i(\delta''_i))}
\end{aligned}\right.
\]

% margin note written sideways
\margin{de c) on conclut par récurrence transfinie\\
$(c^*)\left\{\begin{aligned}\delta'_i&=\alpha_i(\delta')\cup\deleted{\unsure{\text{diagr. de }\delta'_i}}\ \unsure{\text{images dans }C_i\text{ des}|\text{images des}}\ \deleted{\text{diagrammes}}\inserted{\text{cônes proj}}\text{ de type }d'\deleted{\unsure{\text{ contenus}}}\\
\delta''_i&=\alpha_i(\delta'')\cup\deleted{\unsure{\text{diagr. de }\delta''_i}}\ \unsure{\text{images dans }C_i\text{ des cônes ind.}}\text{ de type }d''.
\end{aligned}\right.$}

d) $\forall$ Catégorie $E\in\mathrm{Ob}\,\Delta$, le foncteur $\varphi\longmapsto\varphi\circ\tau_i$ induit une équivalence
\[
\underline{\mathrm{Hom}}_{\delta_{i+1}}(C_{i+1},E)\longrightarrow\underline{\mathrm{Hom}}_{\delta_i}(C_i,E)
\]

\deleted{La} La construction de $K:(C,\delta)- K(C,\delta)$ sera donnée plus bas. On \deleted{\unsure{doit|peut}} \inserted{va} montrer comment elle permet de terminer.

\deleted{De c) on conclut par récurrence transfinie,}

De d), \uline{on déduit par récurrence transfinie}
\[
(d^*)\qquad\boxed{\underline{\mathrm{Hom}}_{\delta_i}(C_i,E)\xrightarrow{\sim}\underline{\mathrm{Hom}}_\delta(C,E)}\qquad \varphi\longmapsto\varphi\circ\alpha_i
\]

Donc on a terminé si on montre qu'il existe $\omega$ \inserted{ordinal} tel que
\[
\left\{\begin{aligned}
1^\circ)&\ C_{\deleted{\unsure{i}}\omega}\in\mathrm{Ob}\,\Delta\qquad 2^\circ)\ \alpha_\omega:C\to C_\omega\ \deleted{\text{est}}\ \delta\text{ exact}\\
3^\circ)&\ \underline{\mathrm{Hom}}_{\delta_\omega}(C_\omega,E)=\underline{\mathrm{Hom}}_{\Delta,\tau_\omega(\delta)}(C_\omega,E)
\end{aligned}\right.
\]

\deleted{\unsure{$2^\circ$) Soit}} Soit pour ceci $\pi$ un cardinal tel que
\[
\pi>\mathrm{card}\,\mathrm{Fl}\,\Theta,\ \deleted{\unsure{\text{card Ob}}}\left\{\begin{aligned}
&\text{si }\Theta\ \deleted{\unsure{\in d'\text{ ou}}}\ \inserted{\in d'\text{ ou }\Theta\in d''}\\
&\text{ou }\Theta\ \deleted{\unsure{\in}}\text{ indexe un }\in\delta'\text{ ou un }\in\delta''
\end{aligned}\right.
\]
et $\omega$ le plus petit ordinal de cardinal $=\pi$. Donc l'ens $S_\omega$ des

% archive sheet number: 4
\scanpage{6}

$j<\omega$ est grand devant $\mathrm{card}\,\mathrm{Fl}\,\Theta$, pour tt $\Theta$ comme dessus. On va établir d'abord \deleted{\unsure{à}} le point suivant, plus fort que $2^\circ$ :

\[
(*)\quad\forall j<\omega,\qquad C_j\unsure{-|\to}C_\omega\text{ est }\delta_j\text{-exact.}
\]

Cela résulte aisément de a) et c) \inserted{(et grandeur de $\omega$)} / \uline{Détaillons la démonstration}

On \deleted{\unsure{puisse montrer|peut}} \inserted{en conclut} $1^\circ$), utilisant $(*)$ et b)

Pour prouver $3^\circ$, on note que \inserted{b) et} c*) [et la grandeur de $\omega$] impliquent
\[
(**)\quad\delta_\omega=\deleted{\unsure{\text{diagr.}}}(\text{cônes exacts de types }d'\text{ et }d'')\cup\alpha_\omega(\delta).
\]
\unsure{$\cup$}

Tout revient donc à construire $K(C,\delta)$, de façon à satisfaire a) : d). On construira \inserted{$(C,\delta)\longmapsto$} $K(C,\delta)$ comme composé $(C,\delta)\longmapsto\sigma\rho(C,\delta)$ de deux constructions, dont chacune fait environ la moitié du boulot :

\inserted{Dans chacune, on doit vérifier l'analogue de c) et d) (qui sera vrai pour le composé).}

$(C,\delta)\longmapsto\rho(C,\delta)$ se fait en plusieurs étapes. On considère $\left\{\begin{aligned}\Delta^{\unsure{\pi}}_{(\varnothing,d'')}&=\Delta''\\\Delta_{(d',\varnothing)}&=\Delta'\end{aligned}\right.$

\deleted{\unsure{$\rho(C,\delta)$ :}}
\[
\begin{aligned}
\rho_1(C,\delta)&=(C_1,\delta_1),\quad\text{avec }C_1\text{ la }\Delta''\text{-enveloppe de }C,\quad\deleted{\unsure{\rho_1}}\inserted{\tau_0}:C\to C_1\\
\delta'_1&=\tau_0(\delta'),\quad\delta''_1=\tau_0(\delta'')\cup(\text{diagrammes de type }d''\text{ de }C_1)
\end{aligned}
\]
\uline{Cela assure} b$'$)\\
c) et d) vérifiés

$\rho_2(C_1,\unsure{d_1|\delta_1})=(C_2,\delta_2)$, \inserted{$\tau_1:C_1-C_2$} avec $C_2$ déduit de $C_1$ en égalisant, tout couple de flèches $z\overset{u}{\underset{v}{\rightrightarrows}}x$ dans $C_1$, de but le sommet $x$ d'un $(x\xrightarrow{p_\lambda}x_\lambda)\in\delta''_1$, satisfaisant $p_\lambda u=p_\lambda v$ pour tt $\lambda$. On \unsure{prend|pose|a} $\delta_2=\tau_1(\delta_1)$

\uline{Cela assure} $a'_1$. c) et d) sont vérifiés.

$\rho_3(C_2,\delta_2)=(C_3,\delta_3)$, \inserted{$\tau_2:C_2-C_3$} avec $C_3$ déduit de $C_2$ en rajoutant, pour $\forall$ \deleted{\unsure{$u$}} \inserted{$\Theta''=$} $(x\overset{p_\lambda}{-}x_\lambda)\in\delta''_2$ et toute \uline{augmentation} $q=(z\xrightarrow{q_\lambda}x_\lambda)_\lambda$ du système projectif $(x_\lambda)$, une flèche $u_{\Theta'',q}$ à $C_3$, satisfaisant $p_\lambda u_{\Theta'',q}=q_\lambda\quad\forall\lambda$.
\[
\delta_3=\tau_2(\delta_2)
\]
\uline{Cela assure} $a'_2$\\
c) et d) vérifiés

% archive sheet number: 5
\scanpage{7}

\[\rho(C,\delta)=\rho_3\rho_2\rho_1(C,\delta)\]

On définit de façon duale une \deleted{opération}\inserted{construction} $\sigma=\sigma_3\sigma_2\sigma_1$, assurant successivement $b'')$, $\deleted{b}a''_1$, $a''_2$.

On pose $K(C,\delta)=\sigma\rho(C,\delta)=\sigma_3\sigma_2\sigma_1\rho_3\rho_2\rho_1(C,\delta)$. Il a les propriétés a) à d) voulues.

\uline{NB} L'ordre dans lequel on a composé les \inserted{6 opérations} $\sigma_i$, \deleted{et les} $\rho_i$ n'a pas d'importance. On pourrait prendre p. ex
\[
K'(C,\delta)=\sigma_1\rho_3\sigma_2\rho_1\rho_2\sigma_3(C,\delta)\qquad !
\]
La chose ayant une signification intrinsèque est la limite des \deleted{\unsure{l'}} itérés transfinis de $K$....

\begin{center}\rule{1.1cm}{0.4pt}\end{center}

\uline{Conditions d'exactitude.}

On a donné \inserted{\unsure{trois paquets|tripliquet|triplet} de $\Delta$-foncteurs entre} des petites catégories dans $\Delta=\Delta_d=\Delta_{d',d''}$, \inserted{(petits)}
\[
f^{(i)}_\lambda:R^{\unsure{(i)}}_\lambda\unsure{-|\to}R^{\prime(i)}_\lambda\quad\inserted{(i=0,1,2)},\quad\text{on pose } f^i=(f_\lambda),\quad f=(f^{(0)},f^{(1)},f^{(2)})
\]

$\Delta_f$ est la sous-$2$-catégorie \inserted{$2$-}pleine de $\Delta$ formée des $E\in\mathrm{Ob}\,\Delta$ tels que, \deleted{le foncteur $\varphi\longmapsto\varphi\vert f$} pour tte $\varphi:R-R'$ dans $\Delta$, les foncteurs $u\longmapsto u\circ\varphi$ de $\underline{\mathrm{Hom}}_\Delta(R',E)\longrightarrow\underline{\mathrm{Hom}}_\Delta(R,E)$. (les flèches entre $E$ \unsure{restant|restent} des foncteurs $d$-exacts)

\[
\begin{aligned}
a_0)&\ \text{sont 0-fidèles si }\varphi\text{ est un } f_\lambda^{\unsure{0}}\\
a_1)&\ \deleted{\text{sont}}\ \unsure{\text{1-fid}|\text{pl-fid}|\text{bijectifs}|\text{injectifs}}\ \deleted{\unsure{\text{si }\varphi\text{ est un}}}\quad f_\lambda^1\\
a_2)&\ \text{'' 2-fid}\quad\text{--- --- --}\quad f_\lambda^2
\end{aligned}
\]

\uline{Théorème} \inserted{Avec les notations précédentes,} Soit $(C,\delta)$ comme dans le th. Alors on peut trouver un foncteur \inserted{$\delta$-exact} $C\overset{i}{-}\bar C$, avec $\bar C\in\mathrm{Ob}\,\Delta_f$, et qui soit \inserted{$2$-}universel pour ces propriétés. De plus, $\bar C$ est ess. petite.

[Cela résoud en particulier les Pb 1 et Pb 2 \inserted{et Pb $2'$} pour $\Delta_f$]

% archive sheet number: 6
\scanpage{8}

\deleted{On peut faire la construction de $K(\bar C,\bar\delta)$ en trois étapes $K(C,\delta)=\rho_2\rho_1\rho_0(C,\delta)$, avec \unsure{sa|par} construction transfinie comme précédemment, en ajoutant les conditions \unsure{$\beta$ e)|$\rho_2$)} aux conditions a) b) c) d)\\ e) \unsure{$e_0$)|$\rho_0$)} Si $R\xrightarrow{\varphi}R'$ est un des $f'_\lambda$, et $u,v:R'\rightrightarrows C_i$.}

Utilisant le th. précédent, on peut supposer que $C\in\mathrm{Ob}\,\Delta$, et $\delta'$ \inserted{$(\delta'')$} contient les \deleted{diagrammes}\inserted{cônes pr.}\inserted{ (ind.)} exacts de type $d'$ ($d''$), de sorte que
\[
\underline{\mathrm{Hom}}_\delta(C,E)=\underline{\mathrm{Hom}}_{\Delta,\delta}(C,E)\qquad\forall E\in\mathrm{Ob}\,\Delta.
\]

On va faire une construction transfinie, $(C_i,\delta_i)$, $C_i$ \inserted{petite} $\in\mathrm{Ob}\,\Delta$ foncteurs de transition dans $\Delta$,
\[
\left\{\begin{aligned}
&(C_0,\delta_0)=(C,\delta)\\
&(C_{i+1},\delta_{i+1})=K(C_i,\delta_i),\quad\tau_i:C_i\unsure{-|\to}C_{i+1},\quad\delta_{i+1}=\tau_i(\delta_i)\\
&C_i=\varinjlim_{j<i\,(\Delta)}C_j\text{ si }i\text{ ordinal limite},\quad\delta_i=\bigcup_{\underline j<i}\text{Images des }\delta_j\text{ dans }C_i.\ \deleted{\unsure{\text{categ.}}}
\end{aligned}\right.
\]

On construit $(C,\delta)\longmapsto K(C,\delta)$ de façon à satisfaire les conditions suivantes

a) \quad $a_0$) Pour tt foncteur $R\xrightarrow{\varphi}R'$ parmi les $f^{(i)}_\lambda$ \inserted{$(i=0,1,2)$}, et pour tout couple de \inserted{$\Delta$-}foncteurs $u',v':R'\rightrightarrows C_i$, et une double flèche $\lambda,\mu:u'-v'$, si $\lambda*\varphi=\mu*\varphi$, alors $\tau_i(\lambda)=\tau_i(\mu)$

$a_1$) Pour tte \deleted{\unsure{$\Delta$- et tte flèche, soit}} $\varphi:R-R'$ parmi les $f^1_\lambda$ et $f^2_\lambda$, et tt couple de $\Delta$-foncteurs $u',v':R'-C_i$, et $\lambda':u'\circ\varphi\to v'\circ\varphi$, \deleted{\unsure{$\exists$ un}} considérant $\tau_i\circ u'\xrightarrow{\tau_i*\lambda'}\tau_i\circ v'$,

\diagram{(\begin{tikzcd}[ampersand replacement=\&,column sep=large] R^{\unsure{\prime|{}}}\arrow[r,shift left=1ex,"\tau_i u'\varphi"{name=U}]\arrow[r,shift right=1ex,swap,"\tau_i v'\varphi"{name=V}]\& C_{i+1}\arrow[from=U,to=V,shorten <=1pt,shorten >=1pt]\end{tikzcd})}

\deleted{\unsure{$\exists$}} $\exists\,\unsure{\bar\lambda|\overset{\text{un}}{\lambda}|\lambda}:\tau_i\circ u\deleted{\circ\varphi}\longrightarrow\tau_i\circ v\deleted{\circ\varphi}$ tel que $\tau_i*\lambda'\deleted{*\varphi}=\deleted{\unsure{\varphi*}}\bar\lambda*\varphi$

$a_2$) Pour tt $\varphi:R\longrightarrow R'$ parmi les $f^2_\lambda$, et tt $\Delta$-foncteur $u:R\longrightarrow C_i$, \deleted{$\exists$} d'où $\tau_i\unsure{:|\circ}u:R-C_{i+1}$, $\exists$ un $\Delta$-foncteur $u':\deleted{\unsure{\Delta}}R'\longrightarrow C_{i+1}$ et un isom
\[
\tau_i\circ u\simeq u'\circ\varphi
\]

% archive sheet number: 7
\scanpage{9}

b) $\tau_i$ est $\delta_i$-exact

c) Pour tout $E\in\mathrm{ob}\,\Delta_f$, le foncteur $\varphi\longmapsto\varphi\circ\tau_i$ de $\underline{\mathrm{Hom}}_\Delta(C_{i+1},E)\longrightarrow\underline{\mathrm{Hom}}_\Delta(C_i,E)$ induit une \uline{équivalence}
\[
\underline{\mathrm{Hom}}_{\Delta,\delta_{i+1}}(C_{i+1},E)\xrightarrow{\approx}\underline{\mathrm{Hom}}_{\Delta,\delta_i}(C_i,E).
\]

\deleted{On pose}

De c) on déduit par récurrence transfinie, pour $E\in\mathrm{ob}\,\Delta_f$
\[
(c^*)\qquad \underline{\mathrm{Hom}}_{\Delta,\delta_i}(C_i,E)\xrightarrow{\approx}\mathrm{Hom}_{\Delta,\delta}(C,E)
\]

et de la définition des $\delta_i$ on déduit par réc. transfinie
\[
\delta_i=\alpha_i(\delta)\qquad(\alpha_i:C\unsure{-|\to}C_i)
\]

Donc \uline{tout revient à prouver que l'on peut trouver un} (petit) \uline{ordinal $\omega$ tel que} $C_\omega\in\mathrm{Ob}\,\Delta_f$, \inserted{\uline{et} $C-C_\omega$ est $\delta$-exact.}

Pour ceci, soit $\pi$ un cardinal majorant strict$^{\mathrm{mt}}$ $\mathrm{Fl}\,R$ et $\mathrm{Fl}\,R'$, \deleted{$\mathrm{Fl}\,\Theta\ \unsure{\mathrm{et}}$} pour les $R,R'$ qui interviennent dans $f$, ainsi que les \inserted{$\mathrm{Fl}\,\Theta$ si $\Theta$ intervient dans $\delta'$ ou $\delta''$.} Soit $\omega$ le plus petit ordinal de cardinal $\pi$. Donc l'ens $S_\omega$ des $j<\omega$ est grand devant $\mathrm{card}\,\mathrm{Fl}\,\unsure{,|{}}R$, $\mathrm{card}\,\mathrm{Fl}\,R'$, $\mathrm{Fl}\,\Theta$.

Mais on voit que $C-C_\omega$ est $\delta$-exact en utilisant b) et le fait que $S_\omega$ grand devant les \inserted{$\mathrm{card}$} $\mathrm{Fl}\,\Theta$. Et on voit que pour $\varphi:R-R'$, $\underline{\mathrm{Hom}}(R',C_\omega)-\underline{\mathrm{Hom}}(R,C_\omega)$ est resp. fidèle ($\varphi$ un $f^{(i)}_\lambda$, avec $i=0,1,2$) ou pl. fidèle $(i=1,2)$, ou équivalence $(\unsure{i|\tau}=2)$, en utilisant respectivement $a_0,a_1,a_2$ plus \unsure{les|le} fait que $S_\omega$ est grand devant \deleted{\unsure{$\mathrm{Fl}$}} $\mathrm{card}\,\mathrm{Fl}\,R$ et $\mathrm{card}\,\mathrm{Fl}\,R'$

% slanted
\margin{dégager lemme sur syst. inductifs de catégories (qui servirait dans les 2 théorèmes...)}

% archive sheet number: 8
\scanpage{10}

Il reste donc à faire la construction de $K(C,\delta)$

% archive sheet number: 9
\scanpage{11}

\uline{Conditions d'exactitude}\par 1) Sur catégories $C$.

a) Existence inconditionnelle de certaines $\varprojlim$ ou $\varinjlim$ \inserted{dans $C$}, correspondant à certains types de diagrammes $d=(d',d'')$.

% slanted
\margin{\uline{Exemple type} existence de $X\unsure{\times|<}Y$ ?}

% slanted
\margin{existence d'objet\unsure{s} \unsure{limites|plats}}

b) Existence conditionnelle de types de $\varprojlim$ ou $\varinjlim$. On se donne \deleted{\unsure{des}}\inserted{un} diagramme $\Theta$, et \deleted{\unsure{des hom et|des sous-}} un paquet de diagrammes \inserted{dans $\Theta$} $\Theta'_i\overset{u'_i}{\unsure{-|\to}}\Theta,\ \Theta''_j\xrightarrow{u''_j}\Theta$, et on demande que pour tt $\Theta\to C$ (satisfaisant éventuellement des conditions supplémentaires de commutativité de certains diagrammes déduits de $\Theta$ -- celà \unsure{se ramène|se revient|revient} à demander que $\Theta$ soit une catégorie et $\Theta\to C$ un foncteur [ voire des conditions ``d'exactitude'' plus subtiles, cf plus bas] conditions d'ex. des foncteurs] les $\varprojlim$ ($\varinjlim$) des diagrammes correspondants de type $\Theta'_i$ ($\Theta''_j$) existent dans $C$.

% slanted
\margin{existence de $X\times_Y X$}

c) Exactitude de cônes projectifs et inductifs. On se donne une \deleted{diagramme}\inserted{catégorie} $\Theta$, muni d'\deleted{\unsure{un}} quatre paquets $\underbrace{\delta',\lambda'}_{\mathrm{c.p.}},\underbrace{\delta'',\lambda''}_{\mathrm{c.i.}}$ de cônes projectifs et inductifs, et on demande que tt foncteur $\delta=(\delta'\unsure{-|,}\delta'')$-exact soit $\lambda=(\lambda',\lambda'')$-exact. [i.e. \inserted{pour} \deleted{un} foncteur \deleted{\illegible{1}}\inserted{$\Theta\to C$}, un certain type d'exactitude en \deleted{\unsure{entraîne}}\inserted{implique} un \deleted{\unsure{en}} autre ]

% slanted
\margin{sommes universelles\\ exactitude des cônes\\ tout morphisme un mono (épi) iso...)}

d) On se donne une \deleted{catégorie}\inserted{foncteur} $\Theta\xrightarrow{\varphi}\bar\Theta$ entre catégories et dans $\Theta$ un paquet $\delta=(\delta',\delta'')$, dans $\bar\Theta$ un paquet $\bar\delta=(\bar\delta',\bar\delta'')$. On considère $u\longmapsto u\circ\varphi$

\diagram{\begin{tikzcd}[ampersand replacement=\&,column sep=large,row sep=0.25em]
\underline{\mathrm{Hom}}(\bar\Theta,C)\arrow[r,"\varphi_C^{*}"]\&\underline{\mathrm{Hom}}(\Theta,C)\\
\cup\&\cup\\
\underline{\mathrm{Hom}}_{\bar\delta}(\bar\Theta,C)\&\underline{\mathrm{Hom}}_{\unsure{\sigma|\delta}}(\Theta,C)
\end{tikzcd}}

% archive sheet number: 10
\scanpage{12}

et on peut demander des conditions des types suivants

(i)\deleted{\unsure{(ii)}} la restriction de $\varphi^*_C$ à $\underline{\mathrm{Hom}}_{\bar\delta}$ est fidèle (cas particulier de c)

(\unsure{$i'$|$ii$})\deleted{\unsure{(iii)}} la restriction de $\varphi^*_C$ à $\underline{\mathrm{Hom}}_{\bar\delta}$ est \deleted{\unsure{un}} surjectif sur les Hom

\inserted{(ii) $\varphi^*_C$ envoie $\underline{\mathrm{Hom}}_{\bar\delta}$ dans $\underline{\mathrm{Hom}}_\delta$ : cas particulier de c), et permet d'exprimer c)}

($ii'$)\deleted{\unsure{(iv)}} $\underline{\mathrm{Hom}}_\delta$ est dans l'image essentielle de $\underline{\mathrm{Hom}}_{\bar\delta}$

\uline{NB} (\unsure{$i'$|$i$}) \deleted{et (ii)} est dans la nature \inserted{d'une condition} d'existence de flèches (ou systèmes de flèches) ds $C$, \deleted{\unsure{et}} ($ii'$) dans la nature d'existence d'objets et de diagrammes d'objets. Ainsi, ($ii'$) généralise a) et b) [ \deleted{\unsure{(ii$'$) impli}} des conditions de type ($i'$) peuvent-elles \inserted{être} exprimées comme conditions de type ($ii'$) ? ]

\margin{? !}

Donc finalement on revient à deux types de conditions permettant d'exprimer les autres

\[
\begin{aligned}1^\circ)&\ \deleted{\unsure{\varphi^*}}\\2^\circ)&\end{aligned}
\left\{\begin{aligned}
&\mathrm{Im\,ess}(\varphi^*_C\vert\underline{\mathrm{Hom}}_{\bar\delta})\subset\underline{\mathrm{Hom}}_\delta\\
&\mathrm{Im\,ess}(\varphi^*_C\vert\underline{\mathrm{Hom}}_{\bar\delta})\supset\underline{\mathrm{Hom}}_\delta
\end{aligned}\right.
\begin{aligned}
&\text{implique conditions d'exactitude }\inserted{\text{habituelles}}\\
&\text{implique }\inserted{\text{conditions d'}}\text{existence (y compris conditionnelle) de limites}
\end{aligned}
\]

Vérifier que cel permet \inserted{bien} d'exprimer l'\unsure{injectivité|intégralité|égalité} resp. surjectivité de $\varphi^*_C\vert\underline{\mathrm{Hom}}_\delta$ sur les Hom ....

$2^\circ$) \uline{Conditions d'exactitude sur foncteurs} $\deleted{\unsure{g}}F:C\longrightarrow C'$

On a déjà vu comment on peut en exprimer, avec des paquets de \deleted{diagrammes}\inserted{cônes} de $C$. Mais on veut des conditions \deleted{\unsure{ind}} d'ex valables pour \uline{tt} $C$, sans donnée supplémentaire sur $C$, -- donc des façons d'associer à tt $C$ des paquets de \deleted{diagrammes}\inserted{cônes} de $C$. Une façon de le faire est de partir de $\Theta$, muni de $\delta\deleted{\unsure{\xi}}=(\delta',\delta'')$ et $\lambda=(\lambda',\lambda'')$ \inserted{(paquets de cônes)} et de prendre les $u:\Theta\unsure{-|\to}C$ qui sont $\delta$-exacts, et demander que $F$ soit $u(\lambda)$-exact pour tt tel $u$\deleted{\unsure{ précis}}. On peut prendre un paquet de \unsure{ce}\deleted{\unsure{$T_\alpha$}} $(\Theta_\alpha,\delta_\alpha,\lambda_\alpha)=T_\alpha$, soit $T=(T_\alpha)$

On désire que \deleted{\unsure{les}} la $T$-exactitude soit propriété

% archive sheet number: 11
\scanpage{13}

par composition \inserted{de foncteurs} entre catégories satisfaisant \unsure{un|aux} ens. également donné de conditions d'exactitude. Pour ceci, il est \deleted{\unsure{bien|bon}} naturel d'exiger que

\diagram{\begin{tikzcd}[ampersand replacement=\&,row sep=1.4em,column sep=1.1em]
\&C\arrow[r,dash,"F"]\&C'\arrow[r,dash,"G"]\&C''\\
\&\Theta\arrow[u,"u"']\&\bar J'\arrow[l,"\eta'"']\&{}\arrow[d,phantom,"{\bigg\}\ \in\lambda=(\lambda',\lambda'')}" description]\\
\bar I'\arrow[ur,"\xi'"]\&\bar I''\arrow[u,"\xi''"']\&\bar J''\arrow[ul,"\eta''"]\&{}\\
{}\arrow[r,phantom,"{\underbrace{\hspace{3.3em}}_{\in\delta=(\delta',\delta'')}}" description]\&{}\&\&
\end{tikzcd}}

\deleted{\unsure{les foncteurs sont les|les foncteurs}} $\lambda'\supset\delta',\lambda''\supset\delta''$. Si on veut que les foncteurs identiques satisfassent ces conditions, il est naturel d'exiger qu'on se borne aux $C$ qui satisfont la condition d'exactitude que tt $\Theta\to C$ $\delta$ exact est $\lambda$ exact. Mais alors la condition \deleted{sur} $F$ devient que $\forall\,\Theta\overset{u}{\unsure{-|\to}} C$ $\lambda$-exact, $F\circ u$ est $\lambda$-exact.

\deleted{On a ainsi donné des moyens de définir}

\deleted{$1^\circ$) \unsure{De} \inserted{un ens $\Sigma_0$} de parties \inserted{$\Delta_0$} de $\mathrm{Ob}(\mathrm{Cat})$, chacune stable par équivalence de catégories (si $C\in\Delta_0$, et $C'\in\mathrm{Ob}\,\mathrm{Cat}$, $C'\cong C$, alors $C'\in\Delta_0$) $\Sigma$ est stable par intersection.}

\deleted{$2^\circ$) Un ens $\Sigma_1$ de parties $\Delta_1$ de $\mathrm{Fl}(\mathrm{Cat})$, chaque $\Delta_1$ stable par isom (si $u,v:C\rightrightarrows C'$, et $u\simeq v$, alors $u\in\Delta_1\Rightarrow v\in\Delta_1$) \uline{contenant les} \unsure{identi|identité} équivalences et par composition}

\textcircled{1} Donnée d'un paquet \inserted{$\rho$} de catégories $\Theta$, avec $\lambda=(\lambda',\lambda'')$. Ici $\rho$ désigne \deleted{\unsure{les}} l'ens des couples $(\Theta,\lambda)$

\deleted{$\Delta_\rho\subset(\mathrm{Cat})$ : les catégories qui}

$\Sigma_\rho\subset\mathrm{Fl}(\mathrm{Cat})$ : un foncteur $F:C-C'$ est $\rho$-exact ssi pour tt \inserted{$(\Theta,\lambda)$ de $\rho$ et tt} foncteur $u:\Theta\to C$ qui est $\lambda$-exact, $F\circ u$ est $\lambda$-exact.

Alors \quad a) $\Sigma_\rho$ est stable par isom (i.e. si $F,G:C-C'$, \deleted{\unsure{et}} $F\in\Sigma_\rho\Rightarrow G\in\unsure{\Sigma|\xi}_\rho$)

b) $\Sigma_\rho$ contient les équivalences de catégories

c) $\Sigma_\rho$ stable par composition

d) Une intersection de parties $\Sigma_\rho$ de $\mathrm{Fl}\,\mathrm{Cat}$ est de $=$ type.

\textcircled{2} Donnée d'un \deleted{\unsure{d'un paquet|de foncteurs}}\inserted{un couple $\sigma\ \deleted{\unsure{\text{de}}}\ (\sigma_1,\sigma_2)$} de paquets de foncteurs $R\xrightarrow{\varphi}\bar R$, avec \deleted{\unsure{$\varphi$}} $R$ muni de $\delta=(\delta',\delta'')$, $\bar R$ de $\bar\delta=(\bar\delta',\bar\delta'')$, et \deleted{\unsure{tels et que}} l'on définit $\Delta_\sigma\subset\mathrm{Ob}\,\mathrm{Cat}$. par

% archive sheet number: 12
\scanpage{14}

$C\in\Delta_\sigma$ ssi

a) si $\unsure{\varphi}\ (R,\delta\,;\bar R,\bar\delta,\varphi:R\unsure{-|\to}\bar R)\in\sigma_1$, alors
\[
\mathrm{Im\,ess}(\varphi_C^*\vert\underline{\mathrm{Hom}}_{\bar\delta}(\bar R,C))\subset\underline{\mathrm{Hom}}_\delta(R,C)
\]

b) si $(-)\in\sigma_2$, alors on a l'inclusion opposée $\supset$.

a) $\Delta_\sigma$ stable par équivalence de catégories

b) Une intersection d'\deleted{catégories}\inserted{ensembles} de la forme $\Delta_\sigma$ est de la \unsure{même|mme} forme.

\textcircled{3} $\Delta_{\sigma,\rho}$ est la \deleted{\unsure{$2$}} sous-$2$-catégorie de $\mathrm{Cat}$ dont les objets sont les éléments de $\Delta_\sigma$, et les foncteurs entre \inserted{ses} objets, les foncteurs $\in\deleted{\Delta}\Sigma_\rho$.

a) $\mathrm{Ob}\,\Delta_{\sigma,\rho}$ stable par équivalence

b) Si $C,C'\in\mathrm{Ob}\,\Delta_{\sigma\rho}$, alors $\underline{\mathrm{Hom}}_{\Delta_{\sigma,\rho}}(C,C')\subset\underline{\mathrm{Hom}}_{\mathrm{Cat}}(C,C')$ est une sous-catégorie str. pleine \deleted{de $C$}

(\,La composition des foncteurs de $\Delta_{\sigma\rho}$ induite par celle de $\mathrm{Cat}$)

c) Toute intersection de catégories de la forme $\Delta_{\sigma,\rho}$ est du $=$ type.

% slanted
\margin{sous-2-catégories $\Delta$ ``distinguées''\\ Si a) les $(R,\delta)$ et $(\bar R,\bar\delta)$ \inserted{provenant de} $\in\sigma$ sont aussi \deleted{des}\inserted{parmi les} $(\unsure{\zeta|\Theta},\lambda)\in\rho$\\ b) Si $(\zeta,\lambda)\in\rho$, alors les $\lambda$-limites existent dans $\Delta$, i.e. parmi les \deleted{\unsure{triples de $\sigma$}} quintuples de $\sigma_2$, il y a \deleted{\unsure{ceux}} (avec $\delta=\varnothing,\bar R$ un cône sur $R$) qui expriment l'existence des $\lambda$ limites.}

\uline{Remarque} Soient $\Delta\subset\Delta'$ de la forme $\Delta_{\sigma,\rho}$ et $\Delta_{\sigma',\rho'}$. Soit $R\xrightarrow{\varphi}\bar R$ dans $\Delta'$, et considérons, \inserted{pour $C\in\mathrm{Ob}\,\Delta$,} $u\longmapsto u\circ\varphi$
\[
\underline{\mathrm{Hom}}_{\Delta'}(\bar R,C)\xrightarrow{\varphi^*_{\Delta',C}=\Psi}\underline{\mathrm{Hom}}_{\Delta'}(R,C)
\]

Considérons l'une des conditions suivantes (sur $C$ variable dans $\mathrm{Ob}\,\Delta$)
\[
\left.\begin{aligned}
&\left.\begin{aligned}a)&\ \psi\text{ est fidèle}\\ b)&\ \psi\text{ est surjectif sur les Hom}\end{aligned}\right\}
\text{ conjonction : }\psi\text{ pl. fidèle}\\
&c)\quad\psi\text{ est ess. surjectif sur objets}
\end{aligned}\right\}\text{ conjonction : }\psi\text{ est une équivalence.}
\]

Soit $\Delta_0\subset\Delta$ la sous $2$-catégorie $1$-pleine définie par l'une déterminée des trois conditions précédentes. Alors \uline{$\Delta_0$ est\deleted{elle} distingué ?}

Soit \deleted{\unsure{$\varphi$|$\Delta$}} $\delta$ $(\bar\delta)$ l'ens des \deleted{diagrammes}\inserted{cônes exacts} de $R$ $(\bar R)$ qui sont image des \deleted{$\Theta\to R$} cônes $\lambda$ \deleted{\unsure{ex}} de $\Theta$ par $\Theta\to R$ $(\Theta-\bar R)$ $\lambda$ exacts \deleted{\unsure{dépendant des conditions de|distinctes des conditions dites}} $(\Theta,\lambda)\in\rho'$. Donc $\underline{\mathrm{Hom}}_{\Delta'}(\bar R,C)=\underline{\mathrm{Hom}}_{\bar\delta}(\bar R,C)$, $\underline{\mathrm{Hom}}_{\Delta'}(R,C)=\underline{\mathrm{Hom}}_\delta(R,C)$

% archive sheet number: 13
\scanpage{15}

\[
\left\{\begin{aligned}
\text{\textcircled{a}}&\quad D=(D',D'')\quad\text{ens de }\deleted{\unsure{\text{petites}}}\text{ catégories}\\
\text{\textcircled{b}}&\quad \sigma=(\sigma_1,\sigma_2),\ \sigma_i\ (i=1,2)\quad\text{ens de }\unsure{\text{quintuples}|\text{quintuplets}|\text{quadruplets}}\ (R,\delta^{\deleted{\unsure{2}}},\bar R,\bar\delta,\varphi:R\unsure{-|\to}\bar R),\\
& R,\bar R\text{ cat, }\varphi\text{ foncteur}\\
&\text{où }\delta=(\delta',\delta'')\ \deleted{\unsure{\text{paquets}|\text{paires}}}\text{ ens de cônes pr. et ind. de } R\text{ \uline{de types }}\underline{\in D}\\
&\phantom{\text{où }}\bar\delta=(\bar\delta',\bar\delta'')\quad\text{---}\qquad \bar R\qquad\text{---}\\
\deleted{\text{\textcircled{c}}}&\quad\deleted{\rho\text{ est ens de }(S,\lambda),\text{ avec }S\text{ cat, }\lambda=(\lambda',\lambda'')\text{ ens de cônes pr et ind de }S\text{ \uline{de type }}\underline{\in D},}\\
&\deleted{\text{et en outre }\rho\text{ contenant au moins les }(R,\delta)\text{ et }(\bar R,\bar\delta)\text{ provenant de }\sigma.}
\end{aligned}\right.
\]

% archive sheet number: 14
\scanpage{16}

\deleted{Or toutes ces assertions découlent facilement de la suivante : tt diagramme \inserted{$\delta$} de type $\Theta\in d''$ dans $D$ provient d'un diagramme $\delta_j$ dans un des $C_j$ $(j<i)$, \uline{et} que l'ens. d'indices des $j<i$ est grand devant $\mathrm{card}\,\mathrm{Fl}\,\Theta$ : on voit alors que la $\varinjlim$ de $\unsure{\overrightarrow{\delta_j}|\delta_j|D}$ dans $C_j$ donne une limite dans $D$ [car les $\varprojlim$ de type $\Theta$ commutent aux $\varinjlim$ grandes devant \inserted{$\mathrm{card}$} $\mathrm{Fl}\,\Theta$)].}

\deleted{b) $\alpha_i:C\unsure{-|\to} C_i$ est $\delta'$-exact, i.e. si $(x\overset{p_\lambda}{-}x_\lambda)$ est dans $\delta'$, alors $(\alpha_i(x)\overset{\alpha_i(p_\lambda)}{-}\alpha_i(x_\lambda))$ est un cône proj. \uline{exact}.}

% slanted
\margin{\uline{\unsure{redonner|redonne}}}

\deleted{Or cela résulte aussitôt de la description des flèches de $\varinjlim_{j<i\,(\mathrm{Cat})} C_i$, de a) et b), et du fait que $i=\{j\mid j<i\}$ est grand devant $\mathrm{card}\,\mathrm{Fl}\,\Theta$ -- de sorte que \unsure{tte|tt} syst \inserted{``compatible''} de flèches $z-\alpha_i(x_\lambda)$ provient d'un syst. analogue dans un $C_j$, $j<i$ \inserted{satisfaisant les compatibilités analogues}, et tte double flèche $u,v:z-\alpha_i(u)$ provient d'une double flèche analogue dans un $C_j$, l'égalité $\alpha_i(p_\lambda)u=\unsure{d_i|\alpha_i}(p_\lambda)v$ pour tt $\lambda$ étant déjà vérifiée.}

\deleted{Cela achève la démonstration.}

% slanted
\margin{\uline{Ce ci résoud à la fois le \deleted{\unsure{seul}} Pb 1 et le Pb 2}}

\deleted{\uline{Th} Soient $d',d''$ des petits ensembles de diagrammes, \deleted{\unsure{et}} $\Delta=\Delta_{(d',d'')}$, $C$ une petite catégorie \inserted{munie d'un ens $(\delta',\delta'')$ de cônes projectifs et inductifs}. Alors \deleted{\unsure{le}} il existe une $\Delta$-enveloppe de $(C,\delta',\delta'')$, \deleted{et celle-ci est ens$^{\mathrm{t}}$ petite.} i.e. un foncteur $u:C-C'$ avec $C'\in\mathrm{Ob}\,\Delta$, $i$ $(\delta',\delta'')$ exact, et universel pour ces propriétés [i.e. $\forall E\in\mathrm{Ob}\,\Delta$, $\varphi\mapsto\varphi\circ i$ induit une équivalence \deleted{$\underline{\mathrm{Hom}}_{\Delta,(\delta',\delta'')}(C$} $\underline{\mathrm{Hom}}_{\Delta,(i(\delta'),i(\delta''))}(C',E)\xrightarrow{\sim}\underline{\mathrm{Hom}}_{\unsure{\Delta,}(\delta',\delta'')}(C,E)$.}

\deleted{\unsure{Soit} Soit $\rho$ le foncteur $\Delta_{(\varnothing,d'')}$-enveloppe,\\ $\sigma$ \hspace{.7em}---\hspace{.7em} $\Delta_{(d',\varnothing)}$-enveloppe.}

\deleted{On construit une suite transfinie de catégories $C_i$ sous $C$, chacune munie d'un ens \inserted{$\delta_i=$} $(\delta'_i,\delta''_i)$ de cônes proj. et inductifs. On pose}

\[
\begin{aligned}
&\deleted{C_0=C,\quad\deleted{\delta'_0=\delta''_0=\varnothing}\quad(\delta'_0,\delta''_0)=(\delta',\delta'')}\\
&\deleted{(C_{i+1},\inserted{\delta_{i+1})}\text{ déduit de }(C_i,\delta_i)\text{ ainsi :}}
\end{aligned}
\]

% archive sheet number: 15
\scanpage{18}

\uline{Topos et structures algébriques}

\margin{$\Delta=$ \unsure{Pseudotopos|Pseudo-topos|bons topos} (conditions a) b) c) de Giraud) foncteurs commutant aux $\varinjlim$, et exacts à g.}

1) Topos et structures algébriques définissables par $\varprojlim$ finies.
\[
\Delta_0=\text{cat. avec }\varprojlim\text{ finies}
\]

$T_0$ th. algébrique rel /$\Delta_0$, \deleted{définie}\inserted{représentée} par une petite $C\in\mathrm{Ob}\,\Delta_0$\\
\deleted{\unsure{Th.}}
\[
T_0(E)\simeq\underline{\mathrm{Hom}}_{\Delta_0}(C,E)\qquad\text{si }E\in\mathrm{Ob}\,\Delta_0
\]

Si $E\in\mathrm{Ob}\,\Delta$, on sait qu'on a
\[
T_0(E)=\underline{\mathrm{Hom}}_{\Delta_0}(C,E]\approxeq\underline{\mathrm{Hom}}_\Delta(\widehat C,E)
\left[\simeq\underline{\mathrm{Hom}}_{\mathrm{top}}(E,\widehat C)^\circ\quad\text{si }E\text{ un topos}\right]
\]
Ainsi \deleted{\unsure{$C$ est la $\Delta_0$}}\inserted{le topos $\widehat C$ est} la \deleted{topos} $\Delta_0$-$\Delta$ enveloppe de $C$.

2) \inserted{$\Delta$-}Enveloppe d'une \inserted{petite} catégorie $I$ (Pb. 1)

On prend la $\Delta_0$-enveloppe $\rho_{\Delta_0}(I)\subset\underline{\mathrm{Hom}}(I,\mathrm{Ens})^\circ$ (sous-cat. pleine des foncteurs ``de présentation finie'' -- sous-catégorie pleine engendrée par $I$ via $\varinjlim$ finies -- on peut le faire en deux crans, rajoutant les $\varinjlim$ finies de diagrammes dans $I$, puis prenant facteurs directs.... puis on prend la $\Delta_0$-$\Delta$-enveloppe de $\rho_{\Delta_0}(I)$, savoir $\widehat{\rho_{\Delta_0}(I)}$.

3) \uline{Sous-théorie d'une $\Delta$-théorie} \inserted{(représentée par un topos)} -- et sous-topos (Pb. 2)

Soit $R$ un topos, $\Sigma$ un ensemble de flèches de $R$. On considère
\[
E\longmapsto\underline{\mathrm{Hom}}_{\Delta,\Sigma}(R,E)\qquad E\in\mathrm{Ob}\,\Delta
\]

On sature $\Sigma$ par ST1, ST2, ST3 (SGA 4 I 9.1.7 c)), d'où $\Sigma'$ et on trouve $R\Sigma'^{-1}$, qui est un \uline{topos} \inserted{$R'$}, et on a
\[
\underline{\mathrm{Hom}}_{\Delta,\Sigma}(R,E)\simeq\underline{\mathrm{Hom}}_\Delta(R',E)
\]

% margin note written sideways
\margin{\uline{ST1} Si $w=vu$, \deleted{\unsure{lorsque}} et \deleted{\unsure{l'une}}\inserted{deux} des flèches $u,v,w$ \deleted{sont} dans $\Sigma$, \deleted{\unsure{alors la}} la troisième l'est. \deleted{\unsure{Contient|Contravariance}} isos}

% margin note written sideways
\margin{\uline{ST2} Stabilité par changement de base (donc par $\varprojlim$ finies)}

% margin note written sideways
\margin{\uline{ST3} Nature locale ($\Leftrightarrow$ stabilité par $\varinjlim$)}

On sait que $R'$ est un sous-topos de $R$, et tout sous-topos de $R$ s'obtient ainsi (loc. cit.). Si $R\simeq\widetilde C$, $C$ petit site,

% archive sheet number: 17
\scanpage{19}

ils correspondent aux topologies \inserted{$T'$} sur $C$ \deleted{\unsure{plus}}\inserted{plus} fines que la top. donnée, par $R'\longleftrightarrow\widetilde{(C,T')}$.

Remarquons que pour un sous topos donné \inserted{$R'$}, on peut tjrs le décrire à l'aide d'un \uline{petit} ens. $\Sigma\subset\mathrm{Fl}\,R$. En effet, soit $C\subset R$ \inserted{petite et} génératrice \deleted{\unsure{de}}. Alors les $\Delta$-flèches $\in\Delta$ $R\longrightarrow\unsure{E|R'|\bar R}$ qui se factorisent par $R'$ sont celles qui transforment familles \deleted{couvrantes} $S_i\to S$ \inserted{dans $C$ couvrantes pour la topologie $T'$ définissant $R'$} \deleted{\unsure{de $R$}} en familles couvrantes [ou encore telles que pour $\forall S\in\mathrm{Ob}\,R$ et $S'\hookrightarrow S$ monomorphisme de \deleted{\unsure{$E$}} $R$, $S'\to S$ est transformé en un \uline{iso}].

\uline{Conséquence} Supposons $R$ ``algébrique'' i.e. $T_R$ \inserted{théorie} algébrique, alors $T_{R'}\subset T_R$ (ss-\deleted{\unsure{catég|catégorie}}théorie ``pleine'') est algébrique, déduite de $T_{R'}$ par un petit ensemble d'axiomes.

D'ailleurs, si on a deux théorie représentables $T_R,T_{R'}$ ($R,R'$ des topos) et un foncteur $T_{R'}\unsure{-|\to} T_R$ qui est un plongement en théories, i.e. tel que \deleted{$\underline{\mathrm{Hom}}_\Delta$} $T_{R'}(E)-T_R(E)$ pl. fid. pour tt $E\in\mathrm{Ob}\,\Delta$, \deleted{\unsure{i.e.} $\underline{\mathrm{Hom}}_\Delta($} alors la flèche correspondante $R\to R'$ de $\Delta$ est telle que $\forall E\in\mathrm{Ob}\,\Delta$, $\underline{\mathrm{Hom}}_\Delta(R',E)-\underline{\mathrm{Hom}}_\Delta(R,E)$ est pl. fidèle. Or ceci caractérise les morphismes de topos $R'-R$ qui sont des plongements.

4) \uline{\unsure{Théorèmes|Théorème}} a) Un $R\in\mathrm{Ob}\,\Delta$ est algébrique, i.e. représente une théorie algébrique, ssi $R$ est un topos, i.e. admet une petite sous-catégorie génératrice. .

b) Soit $I$ un cat discrète, $R\in\mathrm{Ob}\,\Delta$, $I\overset b{-} R$ un foncteur. Alors \unsure{$*$} $b$ est basique (i.e. ...) \deleted{\unsure{désigne|disque}}\deleted{\unsure{ssi} $b(I)$ est génératrice} ssi la sous-catégorie pleine de $R$ engendrée par produits finis d'objets de $b(I)$ est strictement génératrice.

\uline{Dém} \deleted{\unsure{Évidemment b) $\Rightarrow$ a)}} \inserted{Évidemment b) $\Rightarrow$ a, donc} prouvons \deleted{\unsure{d'abord}} b).

\uline{Suffisance} On va construire la théorie alg. \inserted{basique} représentée par $R,b$ en plusieurs pas.

% archive sheet number: 18
